\chapter*{Wykaz ważniejszych oznaczeń i skrótów} % section* - ukrywa numerowanie oraz wyklucza ze spisu tresci
\addcontentsline{toc}{chapter}{Wykaz ważniejszych oznaczeń i skrótów}% % reczne dodanie do spisu tresci
\noindent
\textbf{UX} - (User Experience) doświadczenie użytkownika obejmujące jego odczucia i ocenę korzystania z aplikacji, w tym łatwość obsługi oraz możliwość realizacji zamierzonych celów. \newline
\textbf{UI} - (User Interface) interfejs użytkownika, czyli elementy umożliwiające komunikację z aplikacją, takie jak ekrany, przyciski, pola formularzy i komunikaty. \newline
\textbf{Spektrogram} - reprezentacja czasowo-częstotliwościowa sygnału przedstawiająca zmiany jego zawartości częstotliwościowej w czasie. Natężenie koloru obrazuje poziom poszczególnych składowych częstotliwościowych. \newline
\textbf{Macierz pomyłek} - (Confusion Matrix) tabela zestawiająca rzeczywiste klasy próbek z klasami przewidzianymi przez model. Pozwala określić liczbę poprawnych klasyfikacji oraz błędów dla poszczególnych klas. \newline
\textbf{CNN} - (Convolutional Neural Network) konwolucyjna sieć neuronowa wykorzystywana do analizy danych oraz ekstrakcji cech \cite{o2015introduction}. \newline
\textbf{RNN} - (Recurrent Neural Network) rekurencyjna sieć neuronowa przeznaczona do przetwarzania danych sekwencyjnych oraz analizy zależności czasowych \cite{grossberg2013recurrent}. \newline
\textbf{LSTM} - (Long Short-Term Memory) rodzaj rekurencyjnej sieci neuronowej umożliwiającej modelowanie zależności długoterminowych w danych sekwencyjnych \cite{hochreiter1997long}. \newline
\textbf{GRU} - (Gated Recurrent Unit) uproszczona architektura rekurencyjnej sieci neuronowej wykorzystująca mechanizmy bramkujące do efektywnego przetwarzania sekwencji danych \cite{cho2014learning}. \newline
\textbf{Transformer} - architektura sieci neuronowej oparta na mechanizmie self-attention, stosowana do przetwarzania danych sekwencyjnych oraz analizy zależności długoterminowych \cite{vaswani2017attention}. \newline
\textbf{Flutter} - otwartoźródłowy zestaw narzędzi programistycznych służący do tworzenia natywnych aplikacji mobilnych, webowych oraz desktopowych. \newline
\textbf{Dart} - obiektowy język programowania zoptymalizowany pod kątem tworzenia interfejsów użytkownika oraz aplikacji mobilnych, webowych i desktopowych. \newline
\textbf{iOS} - (iPhone Operating System) mobilny system operacyjny opracowany przez firmę Apple dla urządzeń mobilnych. \newline
\textbf{BEM} - (Breath Extraction Module) moduł ekstrakcji oddechu. \newline
\textbf{BPDM} - (Breathing Phases Detection Module) moduł detekcji faz oddechu. \newline
\textbf{ReLU} - (Rectified Linear Unit) jedna z najczęściej stosowanych funkcji nieliniowych w sieciach neuronowych. Przyjmuje wartość równą zero dla argumentów ujemnych oraz wartość równą argumentowi wejściowemu dla wartości nieujemnych. \newline